\documentclass[11pt]{article}
\usepackage[margin=1in]{geometry}
\usepackage{amsmath,amssymb,amsthm,booktabs,array}
\usepackage[T1]{fontenc}
\usepackage{lmodern}
\usepackage{microtype}
\usepackage[colorlinks=true,allcolors=blue,breaklinks=true]{hyperref}
\newtheorem{theorem}{Theorem}
\newtheorem{proposition}[theorem]{Proposition}
\newtheorem{corollary}[theorem]{Corollary}
\newcommand{\R}{\mathbb R}
\newcommand{\supp}{\operatorname{supp}}
\newcommand{\E}{\mathbb E}
\title{Counterexamples to Ekene E.'s accumulated sparse-error conjectures and a sharp extension of Barron's bound}
\author{Haobo Ma\thanks{ChronoAI Pte Ltd, Singapore.}\and
Wenlin Zhang\thanks{National University of Singapore; ChronoAI Pte Ltd, Singapore.}}
\date{18 November 2026}
\begin{document}
\maketitle
\begin{abstract}
Ekene E.'s finite-dimensional and Hilbert-space conjectures in
\emph{A bound on accumulated PSD contractions}, MathOverflow question 509794,
revision 2, paragraphs 1--2 (2 April 2026)~\cite{ekene}, are false:
we construct equicorrelated counterexamples and give a sharp accumulated
extension of Barron's pointwise approximation bound (1993, Lemma 1)~\cite{barron1993}.
The harmonic upper bound has sharp leading constant one, and the optimal
pointwise squared-error envelope is $1/(1+\sqrt d)^2$.
The same errors are realized by uniform measures on unit vectors, with zero
tails; one $64$-atom example has total strictly between $1.431$ and $1.432$.
A countably supported measure in a separable Hilbert space has a divergent
accumulated squared-error series.
\end{abstract}

\section{The two proposed bounds}
Let $K$ be a real symmetric positive semidefinite matrix with Euclidean
operator norm at most one, and let $x\in\R^m$ satisfy
\begin{equation}\label{eq:hyp}
 x_i\geq0,\qquad \|x\|_2=1,\qquad K_{ii}\leq x_i^2.
\end{equation}
Write
\[
 \tau_d(K,x)^2=
 \min_{z\in\R^m,\ |\supp z|\leq d}(x-z)^\top K(x-z).
\]
The first conjecture in~\cite[paragraph 1]{ekene} proposes
$\sum_{d=1}^{m-1}\tau_d(K,x)^2\leq1$.
Here and below the exponent $2$ is part of the squared-error notation.

For a probability measure $\mu$ on the closed unit ball of a real Hilbert
space, with Bochner barycenter $v=\int a\,d\mu(a)$, put
\[
\sigma_d(v)^2=
 \inf_{\substack{a_1,\ldots,a_d\in\supp\mu\\c_1,\ldots,c_d\in\R}}
 \left\|v-\sum_{j=1}^d c_j a_j\right\|^2.
\]
The second conjecture in~\cite[paragraph 2]{ekene} proposes
$\sum_{d=1}^{\infty}\sigma_d(v)^2\leq1$.
The support is the actual topological support of $\mu$; repetitions and
arbitrary real coefficients are permitted.

The first assertion is introduced by ``Prove that''; the second is the
source's stated ultimate conjecture. Its intervening weaker claim is the
finite-dimensional estimate $\tau_d(K,x)^2\leq1/d$.

The empirical approximation method is due to Maurey; Pisier's 1981 account
uses independent samples in the proof of Lemma 2~\cite{pisier}.
Barron~\cite[Lemma 1, p.~934, (13)]{barron1993} gives the explicit
Hilbert convex-hull bound $(b^2-\|f\|^2)/d$ for the infimum of the squared
error when the atoms have norm at most $b$.
Barron, Cohen, Dahmen and DeVore~\cite[Theorems 2.1--2.2, (2.13), (2.19)]{barron}
give explicit greedy approximation bounds, including the corresponding
variance-form rate. These classical pointwise estimates do not imply
summability. Below, the variable-correlation family determines both the
sharp pointwise envelope and the leading constant of the accumulated
finite-dimensional bound, and a direct-sum construction gives divergence
for one countably supported measure.

\section{Exact minima with variable correlation}
Let $J_m$ be the $m\times m$ all-ones matrix and
$H_m=\sum_{j=1}^m1/j$, with $H_0=0$.
\begin{theorem}\label{thm:main}
For every integer $m\geq1$ and real $0<\rho<1$, put
$a=1-\rho$, $C=a+m\rho$, $D_d=a+d\rho$, and set
\[
 K_{m,\rho}=\frac{aI_m+\rho J_m}{m},\qquad
 x_m=\frac1{\sqrt m}(1,\ldots,1).
\]
These satisfy~\eqref{eq:hyp} and $\|K_{m,\rho}\|_{2\to2}\leq1$.
For every $1\leq d\leq m$,
\begin{equation}\label{eq:error}
 \tau_d(K_{m,\rho},x_m)^2=t_{m,d}(\rho):=
 \frac{aC(m-d)}{m^2D_d}.
\end{equation}
The minimum is attained on any $d$-element support by taking
\begin{equation}\label{eq:optimizer}
 z_i=\frac{C}{\sqrt m D_d}
 \quad\hbox{on that support},\qquad z_i=0\quad\hbox{elsewhere}.
\end{equation}
In $\R^{m+1}$, with orthonormal coordinates $e_0,e_1,\ldots,e_m$, set
\[
 a_i=\sqrt\rho\,e_0+\sqrt a\,e_i,\qquad
 \mu_{m,\rho}=\frac1m\sum_{i=1}^m\delta_{a_i},\qquad
 v_m=\frac1m\sum_{i=1}^m a_i.
\]
Then $\sigma_d(v_m)^2=t_{m,d}(\rho)$ for $1\leq d\leq m$, and
$\sigma_d(v_m)^2=0$ for every $d\geq m$.
Consequently the two accumulated errors are equal to
\begin{equation}\label{eq:total}
 S_m(\rho)=\frac{a(m+b)}{m^2}
 \left((m+b)\sum_{d=1}^m\frac1{d+b}-m\right),
 \qquad b=\frac a\rho.
\end{equation}
For fixed $\rho$, $S_m(\rho)=(1-\rho)\log m+O_\rho(1)$.
In particular, neither proposed bound can be replaced by any universal
finite constant.
\end{theorem}

\begin{proof}
For every $y\in\R^m$,
\[
 y^\top K_{m,\rho}y=\frac{a\sum_i y_i^2+\rho(\sum_i y_i)^2}{m}\geq0.
\]
The eigenvalue on the all-ones direction is $C/m=\rho+a/m\leq1$;
on its orthogonal complement it is $a/m>0$.
Thus $K_{m,\rho}$ is a positive definite contraction.
Also $(K_{m,\rho})_{ii}=1/m=(x_m)_i^2$ and $\|x_m\|_2=1$.

To optimize over all real sparse vectors, put $z=\sqrt m\,c$ and define
$s=\sum_i c_i$, $q=\sum_i c_i^2$. Direct expansion gives
\begin{align}\label{eq:expansion}
 (x_m-z)^\top K_{m,\rho}(x_m-z)
 &=\rho(1-s)^2+a\sum_{i=1}^m(1/m-c_i)^2\notag\\
 &=\rho+\frac am-2\left(\rho+\frac am\right)s+\rho s^2+aq.
\end{align}
If $|\supp c|\leq d$ and $d\geq1$, Cauchy--Schwarz on the actual nonzero
support gives $q\geq s^2/d$; this includes the zero vector.
With $s_*=dC/(mD_d)$, completing the square yields
\begin{equation}\label{eq:completion}
 (x_m-z)^\top K_{m,\rho}(x_m-z)-t_{m,d}(\rho)
 =a\left(q-\frac{s^2}{d}\right)
   +\frac{D_d}{d}(s-s_*)^2\geq0.
\end{equation}
Taking $c_i=C/(mD_d)$ on any $d$ coordinates and zero elsewhere
makes both terms vanish. This proves the global lower bound and attainment,
without restricting signs or allowing only a prescribed support.

Each $a_i$ has norm one and the atoms are distinct. Their finite set is
closed, every atom has mass $1/m$, and the complement has measure zero.
Hence $\supp\mu_{m,\rho}=\{a_1,\ldots,a_m\}$. The finite measure has Bochner
barycenter $v_m$. For arbitrary $c\in\R^m$ the coordinates give
\begin{equation}\label{eq:hilbert}
 \left\|v_m-\sum_{i=1}^m c_i a_i\right\|^2
 =\rho(1-s)^2+a\sum_{i=1}^m(1/m-c_i)^2.
\end{equation}
Any $d$-tuple of support points consolidates, by summing the coefficients
of repeated atoms, to such a vector $c$ with at most $d$ nonzero entries.
Conversely, enumerate any such support and append fixed support points with
coefficient zero to obtain exactly $d$ terms. Thus the two optimizing sets
in~\eqref{eq:expansion} and~\eqref{eq:hilbert} are equal under
$z=\sqrt m\,c$. This proves the Hilbert formula. For $d\geq m$, use all
$m$ atoms with coefficients $1/m$ and then append zero coefficients.
The error is zero, its minimum possible value.

Finally, $t_{m,d}(\rho)=a(m+b)(m-d)/(m^2(d+b))$ and
$(m-d)/(d+b)=(m+b)/(d+b)-1$, giving~\eqref{eq:total}.
For fixed $b>0$, integration of the decreasing function $1/(t+b)$ gives
\[
 \log\frac{m+b+1}{b+1}\leq\sum_{d=1}^m\frac1{d+b}
 \leq\log\frac{m+b}{b}.
\]
The sum is therefore $\log m+O_b(1)$, while $(m+b)/m=1+O_b(m^{-1})$.
Substitution gives the claimed asymptotic.
The Hilbert series is an ordinary convergent series with an identically zero
tail after $d=m-1$. The case $m=1$ has an empty finite sum and total zero.
\end{proof}

\begin{corollary}\label{cor:witness}
Both accumulated-error conjectures fail for $m=64$, $\rho=1/2$.
\end{corollary}
\begin{proof}
At $\rho=1/2$, the errors and accumulated total reduce to
\[
 t_{m,d}(1/2)=\frac{(m+1)(m-d)}{2m^2(d+1)},\qquad
 S_m(1/2)=\frac{m+1}{2m^2}\bigl((m+1)H_m-2m\bigr).
\]
Partitioning $2,\ldots,64$ into the six blocks
$2^{j-1}<k\leq2^j$ gives $H_{64}\geq1+6/2=4$. Therefore
\begin{equation}\label{eq:certificate}
 S_{64}(1/2)\geq\frac{2145}{2048}>1.
\end{equation}
Exact rational summation further gives
$1431/1000<S_{64}(1/2)<1432/1000$.
\end{proof}

\begin{corollary}\label{cor:bridge}
Every finite instance satisfying~\eqref{eq:hyp} is realized by a probability
measure on a finite subset of the Hilbert unit ball, with exactly the same
best $d$-term squared errors for every $d\geq1$.
\end{corollary}
\begin{proof}
Let $B=K^{1/2}$ and retain the indices $i$ with $x_i>0$.
Give $b_i=Be_i/x_i$ mass $x_i^2$. These masses sum to one, and
$\|b_i\|^2=K_{ii}/x_i^2\leq1$.
When $x_i=0$, $\|Be_i\|^2=K_{ii}=0$, so $Be_i=0$.
Thus the barycenter is $\sum_{x_i>0}x_iBe_i=Bx$.
Every sparse vector $z$ gives $Bz=\sum_{x_i>0}(x_i z_i)b_i$;
conversely each support combination is of this form, choosing one index
for each atom and consolidating repetitions. Dropped indices contribute
zero. Zero padding gives exactly $d$ terms in either direction.
The error is $\|B(x-z)\|^2=(x-z)^\top K(x-z)$.
The actual support is the finite set of distinct $b_i$ with positive mass,
including zero when it occurs.
\end{proof}

\begin{corollary}\label{cor:divergent}
There is a Borel probability measure on the unit sphere of a separable real
Hilbert space with $\sum_{d=1}^\infty\sigma_d(v)^2=+\infty$.
\end{corollary}
\begin{proof}
For $n\geq1$ put $w_n=2^{-n}$ and $m_n=2^{n4^n}$.
In the Hilbert direct sum $H=\bigoplus_{n\geq1}\R^{m_n+1}$, place the
$\rho=1/2$ atoms of Theorem~\ref{thm:main} in block $n$, and call the
resulting measure and barycenter $\mu_n,v_n$.
Set $\mu=\sum_{n\geq1}w_n\mu_n$. The weights sum to one;
all atoms have norm one. Distinct atoms in the same block have distance
one, and those in different blocks have distance $\sqrt2$.
Their union is consequently closed, and is exactly $\supp\mu$ because
each atom has positive mass. The Bochner barycenter is
$v=\sum_{n\geq1}w_n v_n$, a norm-absolutely convergent series.

Orthogonal projection onto block $n$ sends any $d$-term approximation
to a combination of at most $d$ atoms of $\mu_n$.
After dividing its coefficients by $w_n$, this gives
$\sigma_d(v)^2\geq w_n^2\sigma_d(v_n)^2$ for every $d\geq1$.
For the half-correlation family,
$S_m(1/2)\geq\tfrac12\log m-2$, directly from its harmonic formula
and $H_m\geq\log m$. Therefore
\[
 \sum_{d=1}^{m_n-1}\sigma_d(v)^2
 \geq w_n^2 S_{m_n}(1/2)
 \geq\frac{n\log2}{2}-\frac2{4^n}\longrightarrow+\infty.
\]
Nonnegativity of the errors proves divergence.
\end{proof}

\section{The pointwise estimate and a corrected bound}
The weaker $1/d$ estimate mentioned in the same source is consistent with
the counterexamples. The following form also gives a dimension-dependent
replacement for the finite accumulated bound.

\begin{proposition}\label{prop:sampling}
Under~\eqref{eq:hyp}, for every integer $d\geq1$,
\begin{equation}\label{eq:finite-upper}
 \tau_d(K,x)^2\leq\frac{\operatorname{tr}K-x^\top Kx}{d}\leq\frac1d.
\end{equation}
Consequently, for $m\geq2$,
$\sum_{d=1}^{m-1}\tau_d(K,x)^2
\leq(\operatorname{tr}K-x^\top Kx)H_{m-1}\leq H_{m-1}$.
For a Borel probability measure on the closed unit ball of a separable real
Hilbert space,
\begin{equation}\label{eq:hilbert-upper}
 \sigma_d(v)^2\leq
 \frac{\int\|a\|^2\,d\mu(a)-\|v\|^2}{d}\leq\frac1d.
\end{equation}
\end{proposition}
\begin{proof}
Choose a random coordinate vector $Y=e_i/x_i$ with probability $x_i^2$
over the indices with $x_i>0$. Then $\E Y=x$.
If $x_i=0$, diagonal domination and positivity give $K_{ii}=0$.
Thus $\E(Y^\top KY)=\sum_{x_i>0}K_{ii}=\operatorname{tr}K$.
For independent copies $Y_1,\ldots,Y_d$, their average $Z$ has at most
$d$ nonzero coordinates, and independence cancels all mixed centered terms:
\[
 \E[(x-Z)^\top K(x-Z)]
 =\frac{\E(Y^\top KY)-x^\top Kx}{d}.
\]
The minimum is at most this finite average of feasible errors.
Finally $0\leq x^\top Kx$ and
$\operatorname{tr}K\leq\sum_i x_i^2=1$, proving~\eqref{eq:finite-upper}.

In the Hilbert case, the support has full measure: in a separable metric
space its complement is a countable union of open null sets.
The identity map is bounded and Bochner integrable. Take independent
$A_1,\ldots,A_d$ with law $\mu$ and average them. Almost surely these
are support points, and
\[
 \E\left\|v-\frac1d\sum_{j=1}^d A_j\right\|^2
 =\frac{\E\|A_1-v\|^2}{d}
 =\frac{\E\|A_1\|^2-\|v\|^2}{d}.
\]
The infimum of feasible errors is at most their expectation.
This proves~\eqref{eq:hilbert-upper}.
\end{proof}

\begin{corollary}\label{cor:pointwise}
For each fixed $d\geq1$, the supremum of $\tau_d(K,x)^2$ over all
admissible finite dimensions and pairs is
$1/(1+\sqrt d)^2$. The same supremum holds for $\sigma_d(v)^2$ over
Borel probabilities on separable Hilbert unit balls.
\end{corollary}
\begin{proof}
Write $t=x^\top Kx$, $V=\operatorname{tr}K-t$ in the matrix case, and
$t=\|v\|^2$, $V=\int\|a\|^2\,d\mu(a)-t$ in the Hilbert case.
In either case $t,V\geq0$ and $t+V\leq1$.
Multiply the random average in Proposition~\ref{prop:sampling} by a real
scalar $\lambda$. Its expected squared error is
$(1-\lambda)^2t+\lambda^2 V/d$.
Minimizing over $\lambda$ gives
\[
 \tau_d^2\ \hbox{or}\ \sigma_d^2
 \leq\frac{tV}{V+dt}\leq\frac{t+V}{(1+\sqrt d)^2}
 \leq\frac1{(1+\sqrt d)^2},
\]
with the first fraction defined as zero when $t=V=0$.
For $t+V>0$, put $u=t/(t+V)$; the middle inequality follows from
\[
 1+(d-1)u-(1+\sqrt d)^2u(1-u)
 =\bigl((1+\sqrt d)u-1\bigr)^2\geq0.
\]
For the reverse inequality, hold $d$ fixed, take
$\rho=1/(1+\sqrt d)$ in Theorem~\ref{thm:main}, and let $m\to\infty$.
Its exact minimum tends to
$\rho(1-\rho)/(1-\rho+d\rho)=1/(1+\sqrt d)^2$.
\end{proof}

For a measure on $M$ distinct support points, its barycenter is a
combination of those $M$ points, so $\sigma_d(v)=0$ for $d\geq M$.
Thus~\eqref{eq:hilbert-upper} also gives the corrected total bound
\[
 \sum_{d=1}^\infty\sigma_d(v)^2
 \leq\left(\int\|a\|^2\,d\mu(a)-\|v\|^2\right)H_{M-1}.
\]

\begin{corollary}\label{cor:sharp}
For $m\geq2$ let $R_m$ be the supremum of
$\sum_{d=1}^{m-1}\tau_d(K,x)^2$ over all admissible $m$-dimensional pairs
in~\eqref{eq:hyp}. Then
\[
 R_m\leq H_{m-1},\qquad \lim_{m\to\infty}\frac{R_m}{\log m}=1.
\]
The same lower examples are finite atomic Hilbert measures. For separable
Hilbert measures every finite prefix is bounded by its variance times
$H_N$, and the leading constant one cannot be decreased uniformly.
\end{corollary}
\begin{proof}
The upper bound is Proposition~\ref{prop:sampling}. For $m\geq3$ set
$k=\lceil\log m\rceil\geq2$ and take $\rho=1/k$ in
Theorem~\ref{thm:main}. Formula~\eqref{eq:total} becomes
\[
 S_m(1/k)=\frac{(k-1)(m+k-1)}{km^2}
 \bigl((m+k-1)(H_{m+k-1}-H_{k-1})-m\bigr).
\]
The integral inequalities used above, now with $b=k-1$, give
\[
 H_{m+k-1}-H_{k-1}=\log m-\log k+O(1).
\]
Since $k/\log m\to1$, $k/m\to0$ and $\log k/\log m\to0$, the displayed
formula gives $S_m(1/k)/\log m\to1$.
Squeezing between $S_m(1/k)$ and $H_{m-1}$ proves the limit.
For Hilbert prefixes, sum~\eqref{eq:hilbert-upper} from $d=1$ to $N$.
In the same examples with $N=m-1$, the variance is
$1-\|v_m\|^2=(1-1/k)(1-1/m)\to1$, and their totals have ratio one to
$\log m$. This also proves sharpness of the variance-normalized coefficient.
\end{proof}

\section{Coverage and exact verification}
The two conjectures and the weaker estimate occur in the same revision of
the posing source. Their coverage is as follows.
\begin{center}
\begin{tabular}{@{}>{\raggedright\arraybackslash}p{0.39\linewidth}>{\raggedright\arraybackslash}p{0.55\linewidth}@{}}
\toprule
Source item or consequence & Mathematical conclusion\\
\midrule
\cite{ekene}, paragraph 1: finite accumulated bound &
Refuted-and-corrected: Corollary~\ref{cor:witness} is a counterexample;
Proposition~\ref{prop:sampling} supplies the harmonic replacement.\\[3pt]
\cite{ekene}, paragraph 2: Hilbert accumulated bound &
Refuted (explicit finite atomic counterexample), Corollary~\ref{cor:witness}.
The finite-support replacement is the variance times $H_{M-1}$.\\[3pt]
\cite{ekene}, paragraph 2: weaker finite $1/d$ claim &
Reproved; prior proof: Barron (1993), Lemma 1~\cite{barron1993},
via Corollary~\ref{cor:bridge}. Proposition~\ref{prop:sampling} retains the
second-moment variance.\\[3pt]
Finite-to-Hilbert realization & Not posed separately; proved in
Corollary~\ref{cor:bridge}.\\[3pt]
Divergent series for one measure & Not posed separately; proved in
Corollary~\ref{cor:divergent}.\\[3pt]
Sharp logarithmic leading constant & Not posed; proved in
Corollary~\ref{cor:sharp}.\\[3pt]
Optimal pointwise constant for each $d$ & Not posed; proved in
Corollary~\ref{cor:pointwise}.\\
\bottomrule
\end{tabular}
\end{center}

The $\rho=1/2$, dimension-$64$ statements have corresponding Lean declarations at a
fixed source version. In the
\href{https://github.com/the-omega-institute/trureturing/blob/88be2a120b67a9d9e724eaae970c8263d0a0e9a4/D5/S3/QuadraticForms/AccumulatedSparseRefutation.lean}{quadratic-form file},
\texttt{quadratic} gives the quadratic identity,
\texttt{sparse\_minimum} the attained all-real minima,
\texttt{accumulation\_gt\_one} the strict total, and
\texttt{result} the negation of the full finite assertion. The namespace is
\texttt{D5.S3.QuadraticForms.\allowbreak AccumulatedSparseRefutation}.

In the
\href{https://github.com/the-omega-institute/trureturing/blob/88be2a120b67a9d9e724eaae970c8263d0a0e9a4/D5/S3/HilbertOperators/Approximation/AccumulatedSupportRefutation.lean}{Hilbert-space file},
\texttt{support\_eq} and \texttt{barycenter\_eq} identify the support and
Bochner integral. The declaration \texttt{error\_bridge} equates quadratic
and Hilbert errors. The declarations \texttt{tuple\_to\_coefficients} and
\texttt{coefficients\_to\_tuple} consolidate repeated atoms and pad sparse
vectors with zero coefficients; \texttt{errorSet\_eq} identifies the two
optimizing sets.
The declarations \texttt{errors\_exact}, \texttt{errors\_zero} and \texttt{ordinary\_sum}
give the minima and ordinary series total; \texttt{result} negates the
full Hilbert assertion. Their namespace is
\texttt{D5.S3.HilbertOperators.\allowbreak Approximation.\allowbreak AccumulatedSupportRefutation}.
The variable-$\rho$, arbitrary-$m$ theorem, the finite-to-Hilbert
realization, the divergent countable example, the variance bounds,
the optimal pointwise constant and
the sharp logarithmic coefficient have written proofs above.
The accompanying \path{verify.py} uses exact rational arithmetic to check
the dimension-$64$ sum, optimizer values, Gram identities and
square-completion identities over stated finite ranges of $m,d,\rho$.

\begin{thebibliography}{9}
\bibitem{ekene}
Ekene E., \emph{A bound on accumulated PSD contractions},
MathOverflow question 509794, revision 2, 2 April 2026, paragraphs 1--2.
\url{https://mathoverflow.net/revisions/509794/2}.
Source text licensed CC BY-SA 4.0.

\bibitem{pisier}
G. Pisier, \emph{Remarques sur un r\'esultat non publi\'e de B. Maurey},
S\'eminaire d'analyse fonctionnelle, \'Ecole Polytechnique, 1980--1981,
expos\'e V, 24 April 1981, pp.~1--12, Lemma 2, pp.~V.2--V.3.
\url{https://www.numdam.org/item/SAF_1980-1981____A5_0/}.

\bibitem{barron1993}
A. R. Barron, \emph{Universal approximation bounds for superpositions of
a sigmoidal function}, IEEE Trans. Inform. Theory \textbf{39} (1993),
no.~3, 930--945, Lemma 1, p.~934, equation (13).
\href{https://doi.org/10.1109/18.256500}{doi:10.1109/18.256500}.

\bibitem{barron}
A. R. Barron, A. Cohen, W. Dahmen and R. A. DeVore,
\emph{Approximation and learning by greedy algorithms},
Ann. Statist. \textbf{36} (2008), no. 1, 64--94.
\href{https://doi.org/10.1214/009053607000000631}{doi:10.1214/009053607000000631}.
\href{https://arxiv.org/abs/0803.1718v1}{arXiv:0803.1718v1},
Theorems 2.1--2.2, equations (2.13), (2.19).
\end{thebibliography}

\medskip
\noindent\small\copyright\ 2026 Haobo Ma and Wenlin Zhang.
This note is licensed under the Creative Commons Attribution 4.0
International license (CC BY 4.0).
\end{document}
